\scanpage{40}\mspage{18}

\deleted{\unsure{Dans le|Donc}} \inserted{Ainsi}

\uline{Corollaire} Soit $M$ un monoïde ayant un
objet minimal. Si \deleted{$M$ n'est pas \unsure{co}} $Y=B_M$
est tel que \deleted{tous} les objets $0$-connexes de
$Y^\wedge$ sont $1$-connexes (ou \unsure{simplement|seulement}, que
leur $H^1(x,\unsure{\mathbb Q|Q})=0$), alors $M$ est cofiltrant.

\medskip\centerline{\rule{1.4cm}{0.4pt}}\medskip

\uline{Théorème} Soit $X$ \inserted{$\in\mathrm{Ob}\,\mathrm{Cat}$} une catégorie pseudo-cofiltrante
(hyp. vérifiée si les objets $0$-connexes \inserted{$F$} de $X^\wedge$
sont $1$-connexes, ou si seulement il existe
un groupe \inserted{abélien} $k\ne0$ tel que \unsure{pour|que|l'on} \deleted{\unsure{ait}} $H^1(F,k)\ne0$).
On suppose que $X$ admet \deleted{des \unsure{obj|objets}} un plus
\deleted{objet} petit objet $e$, et que $X_{/e}$ (qui est
également pseudo-cofiltrant) en admet un également,
ou ce qui revient au $=$, que le monoïde
$M=\mathrm{End}_X(e)$ admet un objet minimal $\varphi$. (NB
ces deux hypothèses de minimalité sont satisfaites
si $\mathrm{Ord}\,X$ et $M=\deleted{\unsure{\lambda|\mathrm{End}}}\mathrm{End}_X(e)$ sont finis, en particulier
si $X$ est finie.) Considérons l'opération
de $\mathbb N$ sur $e$ (identifié à un objet de $X^\wedge$)
définie par $\varphi$, $T_n$ opérant par $\varphi^n$, et distinguons
deux cas :

a) Cette opération \deleted{se} fait ``str. sans pts fixes'', i.e.
$\mathbb N\times e\longrightarrow e\times e$ défini par $(n,u)\longmapsto(\varphi^n u,u)$ est
mono. Il revient au $=$ de dire que pour

% archive sheet number: 39
